\chapter{Testing and Multiple Testing}\label{ch:qm:testing-and-multiple-testing}

A research team tries two hundred variants of a momentum signal and presents the best: its backtest $t$-statistic is 3.1,
a one-sided $p$-value of one in a thousand. Had none of the variants any edge, and had they been independent, the largest of
two hundred $t$-statistics would still exceed 3.1 with probability 17.6\%, about one time in six. The variants are not
independent, which helps; the team also chose the sample period, the universe and the cost model, which does not. This
chapter is about what a test can say: size and power and the lemma that makes the likelihood ratio the right statistic; the
many quiet choices that multiply the tests a researcher runs; the corrections that control the chance of any false discovery,
or their proportion; and the tests built for the search over strategies, which respect the correlation between them.

\section{Tests, size and power}

\begin{definition}[Hypothesis test, null hypothesis, $p$-value]\label{def:qm:testing-and-multiple-testing:test}
A \emph{hypothesis test}\index{hypothesis test} is a rule that, from data $X$, decides whether to reject a \emph{null
hypothesis}\index{null hypothesis} $H_0$ (a statement about the law of $X$, such as ``the strategy's mean return is zero'') in favour
of an alternative $H_1$. It rejects when a statistic $T(X)$ exceeds a critical value. The \emph{$p$-value}\index{p-value} of an
observed $t$ is $\P_{H_0}(T \ge t)$: the probability, under the null, of a statistic at least as extreme as the one seen.
\end{definition}

A $p$-value is not the probability that the null is true, and one minus it is not the probability that the strategy works:
both of those need a prior (chapter~14). It is a statement about the data under one hypothesis.

\begin{definition}[Size, power]\label{def:qm:testing-and-multiple-testing:size}
The \emph{size of a test}\index{size of a test} is the largest probability of rejecting when the null is true (a false
positive, or type I error). The \emph{power of a test}\index{power of a test} at an alternative is the probability of rejecting
when that alternative is true; one minus the power is the type II error.
\end{definition}

\begin{theorem}[Neyman--Pearson lemma]\label{thm:qm:testing-and-multiple-testing:np}
For a simple null (density $f_0$) against a simple alternative ($f_1$), the test that rejects when $f_1(x)/f_0(x) > k$, with $k$
chosen so that its size is $\alpha$, is most powerful: no test of size at most $\alpha$ has higher power.
\end{theorem}

\begin{proof}
Let $\phi^*$ be the \omterm{def:qm:testing-and-multiple-testing:lr}{likelihood-ratio test} and $\phi$ any test of size at most $\alpha$ (both indicator functions of the rejection
region). Pointwise $(\phi^* - \phi)(f_1 - kf_0) \ge 0$: where $\phi^* = 1$, $f_1 > kf_0$; where $\phi^* = 0$, $f_1 \le kf_0$.
Integrating, $\text{power}(\phi^*) - \text{power}(\phi) \ge k(\alpha - \text{size}(\phi)) \ge 0$.
\end{proof}

For the mean of normal returns, the ratio $f_1/f_0$ against any positive mean is increasing in the sample mean, so the
one-sided $t$-test is the most powerful test of ``no edge'' against every positive edge at once. With composite hypotheses and
nuisance parameters, the lemma becomes a recipe.

\begin{definition}[Likelihood-ratio test]\label{def:qm:testing-and-multiple-testing:lr}
The \emph{likelihood-ratio test}\index{likelihood-ratio test} of a null that imposes $r$ restrictions on $\theta$ rejects for large
$\mathrm{LR} = 2\bigl(\ell(\hat\theta) - \ell(\hat\theta_0)\bigr)$, where $\ell$ is the log-likelihood, $\hat\theta$ its maximiser and
$\hat\theta_0$ its maximiser under the restrictions.
\end{definition}

By Wilks' theorem, under the null and regularity conditions (the true parameter inside the parameter space), $\mathrm{LR}$ is
asymptotically $\chi^2_r$. The condition matters in practice: testing that a Hawkes excitation (chapter~7) is zero puts the null
on the boundary $\alpha = 0$, and the limit becomes an equal mixture of a point mass at zero and a $\chi^2_1$, which halves
the $p$-value.

\begin{definition}[Kolmogorov--Smirnov test]\label{def:qm:testing-and-multiple-testing:ks}
The \emph{Kolmogorov--Smirnov test}\index{Kolmogorov--Smirnov test} of the null that $X_1, \dots, X_n$ are iid with a fully
specified continuous cdf $F$ uses $D_n = \sup_x|F_n(x) - F(x)|$, $F_n$ the empirical cdf. Under the null, $\P(\sqrt n\,D_n > \lambda)
\to 2\sum_{k \ge 1}(-1)^{k-1}e^{-2k^2\lambda^2}$, whatever $F$; the 5\% critical value is $\lambda = 1.358$.
\end{definition}

A thousand daily returns drawn from a Student $t$ with four degrees of freedom, scaled to 1\% volatility, give $D_n = 0.0615$
against the normal law with the same volatility: above the critical value $1.358/\sqrt{1\,000} = 0.0429$, a $p$-value of
0.0010. Two cautions come with the test. It is weakest in the tails, where $F_n - F$ is small even when the tails are wrong;
and when the mean and variance are estimated from the same data, $D_n$ is smaller than the tabulated law assumes, so the
critical values must be simulated for the fitted family (the Lilliefors correction).

The sharpest use of size and power in a research group is the question of how long to wait.

\begin{proposition}[Sample size for a Sharpe ratio]\label{prop:qm:testing-and-multiple-testing:years}
With $T$ years of independent returns and an annual \omterm{def:qm:estimation:sharpe}{Sharpe ratio} $\mathrm{SR}$, the one-sided test of $\mathrm{SR} = 0$ at size
$\alpha$ has power approximately $\Phi\bigl(\mathrm{SR}\sqrt T - z_{1-\alpha}\bigr)$, and reaches power $1 - \beta$ after
\[
T = \Bigl(\frac{z_{1-\alpha} + z_{1-\beta}}{\mathrm{SR}}\Bigr)^2 \text{ years}.
\]
\end{proposition}

\begin{proof}
By the \omterm{met:qm:estimation:delta}{delta method} of chapter~11, the annualised estimate is approximately $\mathcal N(\mathrm{SR}, 1/T)$ when the daily ratio is
small. The test rejects when $\widehat{\mathrm{SR}}\sqrt T > z_{1-\alpha}$, which has probability $\Phi(\mathrm{SR}\sqrt T - z_{1-\alpha})$;
setting it to $1 - \beta$ gives $T$.
\end{proof}

At 5\% size and 80\% power, a \omterm{def:qm:estimation:sharpe}{Sharpe ratio} of 2 needs 1.5 years, a \omterm{def:qm:estimation:sharpe}{Sharpe ratio} of 1 needs 6.2 years and a \omterm{def:qm:estimation:sharpe}{Sharpe ratio} of 0.5
needs 24.7 years (\cref{fig:qm:testing-and-multiple-testing:power}). Five years of a genuine \omterm{def:qm:estimation:sharpe}{Sharpe ratio} of 1 are detected 72\% of the time: more
than one in four good strategies fails its first honest test.

\begin{omfigure}
\begin{tikzpicture}
\begin{axis}[width=10cm,height=5cm,xlabel={years of daily returns},ylabel={power at 5\% size},
  tick label style={font=\small},label style={font=\small},xmin=0,xmax=15,ymin=0,ymax=1,grid=major,grid style={gray!25},
  legend columns=4,legend style={font=\footnotesize,at={(0.5,-0.32)},anchor=north,draw=none,fill=none,
  /tikz/every even column/.append style={column sep=8pt}},table/col sep=comma]
\addplot[omDef,very thick] table[x=years,y=sr2]{figdata/methods/12-testing-and-multiple-testing/power.csv};
\addplot[omThm,very thick,dashed] table[x=years,y=sr1]{figdata/methods/12-testing-and-multiple-testing/power.csv};
\addplot[black!70,very thick,dotted] table[x=years,y=sr05]{figdata/methods/12-testing-and-multiple-testing/power.csv};
\addplot[black!50,thin,sharp plot] coordinates {(0,0.8) (15,0.8)};
\legend{$\mathrm{SR} = 2$,$\mathrm{SR} = 1$,$\mathrm{SR} = 0.5$,80\% power}
\end{axis}
\end{tikzpicture}

\omcaption{Power of the one-sided 5\% test of a zero \omterm{def:qm:estimation:sharpe}{Sharpe ratio} against years of independent daily returns, for true annual
\omterm{def:qm:estimation:sharpe}{Sharpe ratios} of 2, 1 and 0.5. The line at 80\% power is crossed after 1.5, 6.2 and 24.7 years. Data: the chapter's
tutorial module.}\label{fig:qm:testing-and-multiple-testing:power}
\end{omfigure}

\section{The garden of forking paths}

\begin{definition}[Multiple testing, data snooping, garden of forking paths]\label{def:qm:testing-and-multiple-testing:snooping}
\emph{Multiple testing}\index{multiple testing} is running many tests and reporting some of them. \emph{Data
snooping}\index{data snooping} is choosing a model, a strategy or a parameter by searching over the same data that then test
it, so that the reported statistic is the best of a search presented as a single test. The \emph{garden of forking
paths}\index{garden of forking paths} is the implicit version: every choice the analyst would have made differently had the data
been different (the sample period, the universe, the filter, the cost model) multiplies the tests actually run, even when only
one analysis is ever performed.
\end{definition}

The name is Gelman and Loken's (2013), and the point is theirs: a researcher who never runs two backtests can still snoop,
because the one backtest was shaped by looking at the data. A search is only corrected for if it is counted.

\begin{proposition}[The maximum of independent statistics]\label{prop:qm:testing-and-multiple-testing:max}
If $Z_1, \dots, Z_N$ are independent standard normal, $\P(\max_i Z_i \ge c) = 1 - \Phi(c)^N$, and $\E[\max_i Z_i] \sim
\sqrt{2\ln N}$ as $N \to \infty$.
\end{proposition}

\begin{proof}
$\P(\max_i Z_i < c) = \prod_i\P(Z_i < c) = \Phi(c)^N$. For the growth, $1 - \Phi(c) \approx \varphi(c)/c$, and the maximum sits where
$N(1 - \Phi(c)) \approx 1$, that is $c^2/2 \approx \ln N$ to leading order.
\end{proof}

The asymptote is slow: for $N = 200$ the expected maximum is 2.75, not $\sqrt{2\ln200} = 3.26$. That maximum of two hundred
noise statistics clears the conventional $t > 2$ bar almost always, and the hook's 3.1 about one time in six
(\cref{fig:qm:testing-and-multiple-testing:maxt}). A single test at $t = 3.1$ has $p = 0.001$; the best of two hundred independent tests at $t = 3.1$
has $p = 0.176$.

\begin{omfigure}
\begin{tikzpicture}
\begin{semilogyaxis}[width=10cm,height=5.2cm,xlabel={threshold $c$},ylabel={$\P(\max t \ge c)$ under the null},
  tick label style={font=\small},label style={font=\small},xmin=0,xmax=4.5,ymin=1e-5,ymax=1.5,grid=major,grid style={gray!25},
  legend columns=3,legend style={font=\footnotesize,at={(0.5,-0.32)},anchor=north,draw=none,fill=none,
  /tikz/every even column/.append style={column sep=8pt}},table/col sep=comma]
\addplot[black!70,very thick,dotted] table[x=c,y=single]{figdata/methods/12-testing-and-multiple-testing/maxt.csv};
\addplot[omDef,very thick] table[x=c,y=correlated]{figdata/methods/12-testing-and-multiple-testing/maxt.csv};
\addplot[omThm,very thick,dashed] table[x=c,y=independent]{figdata/methods/12-testing-and-multiple-testing/maxt.csv};
\addplot[black!50,thin,sharp plot] coordinates {(3.11,1e-5) (3.11,1.5)};
\legend{one test,{200 correlated variants},{200 independent tests}}
\end{semilogyaxis}
\end{tikzpicture}

\omcaption{Probability that the largest $t$-statistic reaches $c$ when nothing works: one test, the two hundred momentum
variants of the tutorial (Gaussian simulation with their estimated correlation), and two hundred independent tests. At the
best variant's 3.11 (vertical line): 0.0009, 0.030 and 0.17. Data: the chapter's tutorial, seeded.}\label{fig:qm:testing-and-multiple-testing:maxt}
\end{omfigure}

Real searches are correlated. The tutorial's family is two hundred trend rules that hold the sign of the trailing $L$-day
return, for lookbacks $L = 5, \dots, 204$ days, applied to ten years of a driftless random walk with 1\% daily volatility: no rule
has any edge. Neighbouring lookbacks share almost all their signals (the best variant's daily P\&L has correlation 0.95 with its
neighbour's), distant ones much less (0.28 with the 5-day rule, 0.20 with the 204-day rule). In the history shown, 33 of the
200 variants have a one-sided $p$-value below 5\%, the mean $t$-statistic across the family is 1.04 rather than zero, since
every variant sees the same trending episodes, and the best, at a 42-day lookback, has $t = 3.11$ and an annualised \omterm{def:qm:estimation:sharpe}{Sharpe ratio} of
0.98 (\cref{fig:qm:testing-and-multiple-testing:variants}). That history is the 336th we simulated: we drew histories until the best variant came
near the hook's 3.1, which is the chapter's subject enacted on the chapter itself.

\begin{omfigure}
\begin{tikzpicture}
\begin{axis}[width=11cm,height=5cm,xlabel={lookback $L$ (days)},ylabel={$t$-statistic},
  tick label style={font=\small},label style={font=\small},xmin=0,xmax=210,ymin=-1.5,ymax=4,grid=major,grid style={gray!25},
  legend columns=4,legend style={font=\footnotesize,at={(0.5,-0.32)},anchor=north,draw=none,fill=none,
  /tikz/every even column/.append style={column sep=6pt}},table/col sep=comma]
\addplot[omDef,thick,mark=*,mark size=0.8pt] table[x=lookback,y=t]{figdata/methods/12-testing-and-multiple-testing/variants.csv};
\addplot[black!60,thin,dashed,sharp plot] coordinates {(0,1.645) (210,1.645)};
\addplot[omThm,very thick,dashdotted,sharp plot] coordinates {(0,2.916) (210,2.916)};
\addplot[black,thick,dotted,sharp plot] coordinates {(0,3.481) (210,3.481)};
\legend{variants,one test 5\%,max-$t$ 5\%,Bonferroni 5\%}
\end{axis}
\end{tikzpicture}

\omcaption{The $t$-statistics of the two hundred trend rules on one simulated ten-year random walk, against their lookback, with
three one-sided 5\% thresholds: a single test (1.645), the max-statistic critical value for this correlated family (2.92) and
Bonferroni's (3.48). The best rule, at 42 days, has $t = 3.11$. Data: the chapter's tutorial, seeded.}\label{fig:qm:testing-and-multiple-testing:variants}
\end{omfigure}

\section{Family-wise error and false discoveries}

\begin{definition}[Family-wise error rate, Bonferroni correction, Holm procedure]\label{def:qm:testing-and-multiple-testing:fwer}
For $m$ hypotheses of which $m_0$ are true, the \emph{family-wise error rate}\index{family-wise error rate} (FWER) of a procedure is
the probability that it rejects at least one true null. The \emph{Bonferroni correction}\index{Bonferroni correction} rejects
$H_i$ when $p_i \le \alpha/m$. The \emph{Holm procedure}\index{Holm procedure} orders the $p$-values $p_{(1)} \le \dots \le p_{(m)}$ and
rejects $H_{(1)}, \dots, H_{(k)}$, where $k$ is the last index with $p_{(j)} \le \alpha/(m - j + 1)$ for every $j \le k$.
\end{definition}

\begin{proposition}[Bonferroni and Holm control the family-wise error rate]\label{prop:qm:testing-and-multiple-testing:fwer}
Under any dependence between the tests, both procedures have FWER at most $\alpha$; Holm rejects everything Bonferroni rejects.
\end{proposition}

\begin{proof}
Bonferroni: the union bound over the true nulls gives $\P(\exists i \in \mathcal H_0: p_i \le \alpha/m) \le m_0\alpha/m \le \alpha$.
Holm: let $j^*$ be the rank of the smallest $p$-value among the true nulls. At most $m - m_0$ hypotheses precede it, so $j^* \le
m - m_0 + 1$, and Holm rejects a true null only if $p_{(j^*)} \le \alpha/(m - j^* + 1) \le \alpha/m_0$, which by the union bound over
the $m_0$ true nulls has probability at most $\alpha$. Holm's thresholds are never below $\alpha/m$.
\end{proof}

Controlling the chance of any false discovery is the right goal when one false discovery is expensive, as when a single strategy
goes live. When a research group screens hundreds of signals for a combined model, a few false ones among many true ones cost
little, and demanding none throws away most of the true ones.

\begin{definition}[False discovery rate, Benjamini--Hochberg procedure]\label{def:qm:testing-and-multiple-testing:fdr}
If a procedure makes $R$ rejections of which $V$ are true nulls, its \emph{false discovery rate}\index{false discovery rate} is
$\E[V/\max(R, 1)]$. The \emph{Benjamini--Hochberg procedure}\index{Benjamini--Hochberg procedure} at level $q$ rejects $H_{(1)},
\dots, H_{(k)}$ for the largest $k$ with $p_{(k)} \le kq/m$.
\end{definition}

\begin{theorem}[Benjamini--Hochberg]\label{thm:qm:testing-and-multiple-testing:bh}
If the $p$-values of the true nulls are independent of each other and of the others, the \omterm{def:qm:testing-and-multiple-testing:fdr}{Benjamini--Hochberg procedure} has \omterm{def:qm:testing-and-multiple-testing:fdr}{false
discovery rate} exactly $q\,m_0/m \le q$. It remains at most $q$ under positive regression dependence, and at most $q$ under any
dependence if $q$ is replaced by $q/\sum_{k=1}^m k^{-1}$ (Benjamini and Yekutieli, 2001).
\end{theorem}

\begin{example}[Five $p$-values]\label{ex:qm:testing-and-multiple-testing:five}
Take $p = (0.005, 0.01, 0.03, 0.04, 0.2)$ and $\alpha = q = 5\%$. Bonferroni's threshold $0.01$ rejects two. Holm compares them in
order with $0.05/5, 0.05/4, 0.05/3, \dots$: $0.005 \le 0.01$ and $0.01 \le 0.0125$, then $0.03 > 0.0167$ stops it at two. Benjamini--Hochberg
compares with $0.01, 0.02, 0.03, 0.04, 0.05$: the largest $k$ with $p_{(k)} \le 0.01k$ is $k = 4$, so it rejects four. As adjusted
$p$-values (the smallest level at which each is rejected): Bonferroni $(0.025, 0.05, 0.15, 0.2, 1)$, Holm $(0.025, 0.04, 0.09, 0.09, 0.2)$,
Benjamini--Hochberg $(0.025, 0.025, 0.05, 0.05, 0.2)$.
\end{example}

The trade shows at scale. Among a thousand candidate signals of which a hundred are real (with $t$-statistics centred on 3), the
naive 5\% test makes 136 discoveries of which 45 are false; Bonferroni and Holm make about 18, with a false one in 5\% of
screens; Benjamini--Hochberg at 5\% makes 63, of which 2.9 are false on average, a \omterm{def:qm:testing-and-multiple-testing:fdr}{false discovery rate} of 4.4\% against the
theorem's $0.05 \times 900/1\,000 = 4.5\%$ (\cref{fig:qm:testing-and-multiple-testing:fdr}). It pays for its power with a \omterm{def:qm:testing-and-multiple-testing:fwer}{family-wise error rate} of 92\%:
almost every screen contains a false signal, and that is the design.

\begin{omfigure}
\begin{tikzpicture}
\begin{axis}[width=10cm,height=4.8cm,xbar stacked,bar width=10pt,
  xlabel={discoveries per screen (average of 500)},
  ytick={0,1,2,3},yticklabels={no correction,Bonferroni,Holm,Benjamini--Hochberg},
  tick label style={font=\small},label style={font=\small},
  xmin=0,xmax=150,ymin=-0.6,ymax=3.6,grid=major,grid style={gray!25},
  legend columns=2,legend style={font=\footnotesize,at={(0.5,-0.4)},anchor=north,draw=none,fill=none,
  /tikz/every even column/.append style={column sep=8pt}},table/col sep=comma]
\addplot[fill=omDef!60,draw=omDef] table[x=true,y=k]{figdata/methods/12-testing-and-multiple-testing/fdr.csv};
\addplot[fill=omThm!50,draw=omThm] table[x=false,y=k]{figdata/methods/12-testing-and-multiple-testing/fdr.csv};
\legend{true discoveries,false discoveries}
\end{axis}
\end{tikzpicture}

\omcaption{A thousand signals, a hundred of them real: average true and false discoveries of four rules at 5\%. No correction:
91 true, 45 false. Bonferroni and Holm: 18 true, 0.05 false. Benjamini--Hochberg: 61 true, 2.9 false (a \omterm{def:qm:testing-and-multiple-testing:fdr}{false discovery rate}
of 4.4\%). Data: the chapter's tutorial, seeded.}\label{fig:qm:testing-and-multiple-testing:fdr}
\end{omfigure}

On the momentum family neither control is kind. Bonferroni and Holm both give the best variant an adjusted $p$-value of $200
\times 0.00094 = 0.187$; Benjamini--Hochberg gives 0.076, smaller because its step-up borrows strength from the 33 other small
$p$-values. No variant survives either at 5\%. But both corrections were built for the worst case of dependence, and two hundred
variants that share most of their signals are far from two hundred independent chances.

\section{Reality-check tests for strategy search}

\begin{definition}[Reality Check, superior predictive ability test]\label{def:qm:testing-and-multiple-testing:rc}
For $K$ strategies with performance series $d_{k,t}$ in excess of a benchmark, the \emph{Reality Check}\index{Reality Check}
(White, 2000) tests $H_0: \max_k\E[d_k] \le 0$ with the statistic $\max_k\sqrt n\,\bar d_k$, whose null distribution is estimated by
the bootstrap of $\max_k\sqrt n(\bar d_k^* - \bar d_k)$. The \emph{superior predictive ability test}\index{superior predictive ability
test} (Hansen, 2005) studentises each $\bar d_k$ and recentres only the strategies that are not clearly worse than the benchmark, so
that adding poor strategies to the search does not make the test more conservative.
\end{definition}

Both tests replace the union bound by the actual law of the maximum, and the bootstrap (chapter~13) keeps the dependence between
strategies and over time. When the statistics are approximately jointly normal, a Gaussian simulation with their estimated
correlation does the same job at no cost.

\begin{method}[Max-statistic $p$-values]\label{met:qm:testing-and-multiple-testing:maxt}
Given $t$-statistics $t_1, \dots, t_m$ and the correlation matrix $R$ of the underlying P\&L series: draw $Z^{(s)} \sim \mathcal N(0, R)$
for $s = 1, \dots, S$; the adjusted $p$-value of $H_i$ is the share of draws with $\max_j Z_j^{(s)} \ge t_i$. The step-down version
(Romano and Wolf, 2005) takes the maximum, for the $k$-th largest statistic, only over the hypotheses not yet rejected. Both
control the \omterm{def:qm:testing-and-multiple-testing:fwer}{family-wise error rate}; the effective number of independent trials of the search is the $N$ that solves $1 - (1 -
p_{\text{naive}})^N = p_{\text{adjusted}}$.
\end{method}

For the momentum family, 200\,000 draws from the estimated correlation give the best variant an adjusted $p$-value of 0.030,
against 0.187 for Bonferroni and 0.171 for two hundred independent tests. The search over two hundred correlated lookbacks amounts
to 32 independent trials. The max-statistic 5\% critical value is 2.92, well below Bonferroni's 3.48, and three variants clear it.
The check that matters: over 2\,000 fresh ten-year histories, the best variant reaches $t = 3.11$ in 2.6\% of them, close to the
simulated 3.0\%.

So the corrected test rejects, at 5\%, a strategy that has no edge. It is allowed to, 3\% of the time; and this history was
the 336th we drew. The correction accounts for the two hundred lookbacks it was told about and for nothing else: not the
histories we discarded, not the other signal families the team tried last year. Harvey, Liu and Zhu (2016) reach a similar
conclusion for published equity factors, arguing that a new factor should clear a $t$-statistic of about 3 rather than 2.

\begin{definition}[Deflated Sharpe ratio]\label{def:qm:testing-and-multiple-testing:dsr}
For a \omterm{def:qm:estimation:sharpe}{Sharpe ratio} $\widehat{\mathrm{SR}}$ (per period) estimated from $n$ returns with skewness $\gamma_3$ and kurtosis $\gamma_4$, the
probability that the true ratio exceeds $\mathrm{SR}_0$ is approximately
\[
\Phi\Bigl((\widehat{\mathrm{SR}} - \mathrm{SR}_0)\sqrt{n - 1}\Big/\sqrt{1 - \gamma_3\widehat{\mathrm{SR}} + \tfrac{\gamma_4 - 1}4\widehat{\mathrm{SR}}^2}\Bigr).
\]
The \emph{deflated Sharpe ratio}\index{deflated Sharpe ratio} (Bailey and López de Prado, 2014) is this probability with $\mathrm{SR}_0$
the expected maximum of $N$ null \omterm{def:qm:estimation:sharpe}{Sharpe ratios} of variance $V$: $\mathrm{SR}_0 = \sqrt V\bigl((1 - \gamma_E)\Phi^{-1}(1 - 1/N) +
\gamma_E\Phi^{-1}(1 - 1/(Ne))\bigr)$, $\gamma_E$ the Euler--Mascheroni constant.
\end{definition}

For the best momentum variant the probability that its true \omterm{def:qm:estimation:sharpe}{Sharpe ratio} is positive is 0.999, the single-test view. Deflated for
two hundred trials (with $V = 1/n$, the null variance), the benchmark becomes an annual 0.87 and the \omterm{def:qm:testing-and-multiple-testing:dsr}{deflated Sharpe ratio} 0.63;
for the 32 effective trials, 0.67 and 0.84. Neither reaches the 0.95 a desk would want, and the number of trials, which the
researcher controls and the reader rarely sees, moves the answer more than the data do.

\section{Tutorial: the best of two hundred}\label{tut:qm:testing-and-multiple-testing:best}

\begin{tutorial}
\textbf{Goal.} Search two hundred worthless trend rules, pick the best, and measure how surprising it is under five corrections.
\textbf{End state:} \cref{fig:qm:testing-and-multiple-testing:variants,fig:qm:testing-and-multiple-testing:maxt} and the adjusted $p$-values 0.187, 0.076 and 0.030.
\begin{tutsteps}
\item \textbf{Adjusted $p$-values} by Holm's step-down and Benjamini--Hochberg's step-up.
\omcode{code/firm/multitest/firm_multitest.py}{68}{88}{Holm and Benjamini--Hochberg adjusted $p$-values.}\label{lst:qm:testing-and-multiple-testing:holm}
\item \textbf{Max-statistic $p$-values} from a Gaussian null with the estimated correlation, or from any matrix of null draws.
\omcode{code/firm/multitest/firm_multitest.py}{97}{133}{Single-step and step-down max-statistic $p$-values.}\label{lst:qm:testing-and-multiple-testing:maxt}
\item \textbf{Run} \texttt{problem()} in \texttt{qm\_testing.py}, then \texttt{family\_exceedance(3.11)} to check the simulated $p$-value on
fresh histories, and \texttt{fig\_testing.py}.
\end{tutsteps}
\textbf{What to change next.} Add the mirror-image reversal rules and see the effective number of trials grow; replace the Gaussian
null by the stationary bootstrap of chapter~13 through the \texttt{null\_draws} hook.
\end{tutorial}

\section{Build: the multiple-testing module}\label{bld:qm:testing-and-multiple-testing:multitest}

\begin{build}
\bfield{Purpose} Every strategy the miniature firm promotes from research carries the size of the search behind it and an adjusted
$p$-value; this module computes them.
\bfield{Interface} \texttt{bonferroni(p)}, \texttt{holm(p)}, \texttt{benjamini\_hochberg(p)}, \texttt{benjamini\_yekutieli(p)} returning
adjusted $p$-values; \texttt{maxt\_pvalues(t, corr, n\_sim, seed, null\_draws, stepdown)}; \texttt{effective\_trials(p\_single,
p\_family)}; \texttt{expected\_max\_sr(n\_trials, sr\_var)}; \texttt{deflated\_sharpe(sr, n\_obs, skew, kurt, sr0)}; \texttt{norm\_cdf},
\texttt{norm\_ppf}.
\bfield{Rules} Adjusted $p$-values, not reject flags, so that the level is chosen by the caller; one-sided statistics for strategy
search; the null draws are a parameter, so that the bootstrap replaces the Gaussian without touching the callers.
\bfield{Acceptance tests} \texttt{code/firm/multitest/tests/}: the five-$p$-value example; FWER and FDR control by simulation; the
max-statistic $p$-value equals the independent formula for $R = I$ and the naive $p$-value for perfectly correlated tests; the
expected-maximum approximation within 1.5\%.
\bfield{Stretch} Hansen's SPA test with the stationary bootstrap; the Romano--Wolf step-down with bootstrap null draws; the
probability of backtest overfitting by combinatorial cross-validation.
\end{build}

\begin{omsources}
\item J.~Neyman and E.~S.~Pearson, ``On the problem of the most efficient tests of statistical hypotheses'', \emph{Philosophical
Transactions of the Royal Society A} 231, 1933.
\item S.~Holm, ``A simple sequentially rejective multiple test procedure'', \emph{Scandinavian Journal of Statistics} 6, 1979.
\item Y.~Benjamini and Y.~Hochberg, ``Controlling the false discovery rate: a practical and powerful approach to multiple testing'',
\emph{Journal of the Royal Statistical Society B} 57, 1995.
\item Y.~Benjamini and D.~Yekutieli, ``The control of the false discovery rate in multiple testing under dependency'', \emph{Annals of
Statistics} 29, 2001.
\item H.~White, ``A reality check for data snooping'', \emph{Econometrica} 68, 2000.
\item P.~R.~Hansen, ``A test for superior predictive ability'', \emph{Journal of Business and Economic Statistics} 23, 2005.
\item J.~P.~Romano and M.~Wolf, ``Stepwise multiple testing as formalized data snooping'', \emph{Econometrica} 73, 2005.
\item C.~R.~Harvey, Y.~Liu and H.~Zhu, ``\dots and the cross-section of expected returns'', \emph{Review of Financial Studies} 29, 2016.
\item D.~H.~Bailey and M.~López de Prado, ``The deflated Sharpe ratio'', \emph{Journal of Portfolio Management} 40(5), 2014.
\item A.~Gelman and E.~Loken, ``The garden of forking paths'', working paper, 2013.
\end{omsources}

\section{Exercises}

\begin{exercise}[$\star$]\label{exo:qm:testing-and-multiple-testing:1}
A backtest has $t = 2.5$. What is its one-sided $p$-value, and its Bonferroni-adjusted $p$-value if it is the best of twenty
variants?
\end{exercise}

\begin{exercise}[$\star$]\label{exo:qm:testing-and-multiple-testing:2}
A test rejects when $Z > 1.645$. Under the alternative of interest $Z \sim \mathcal N(2, 1)$. What are its size and its power?
\end{exercise}

\begin{exercise}[$\star$]\label{exo:qm:testing-and-multiple-testing:3}
Apply Bonferroni and Holm at 5\% to $p = (0.001, 0.012, 0.02, 0.3)$.
\end{exercise}

\begin{exercise}[$\star\star$]\label{exo:qm:testing-and-multiple-testing:4}
How many years of independent returns does a test of 5\% size need to detect an annual \omterm{def:qm:estimation:sharpe}{Sharpe ratio} of 0.7 with 90\% power?
\end{exercise}

\begin{exercise}[$\star\star$]\label{exo:qm:testing-and-multiple-testing:5}
Daily returns are normal with mean zero. The first 500 days have sample variance $(1.0\%)^2$, the next 500 $(1.2\%)^2$. Compute the
likelihood-ratio statistic for a common variance and its $p$-value.
\end{exercise}

\begin{exercise}[$\star\star$]\label{exo:qm:testing-and-multiple-testing:6}
For ten independent noise strategies, what is the probability that the best has $t \ge 2$, and what is the expected maximum by the
formula of the \omterm{def:qm:testing-and-multiple-testing:dsr}{deflated Sharpe ratio}?
\end{exercise}

\begin{exercise}[$\star\star\star$]\label{exo:qm:testing-and-multiple-testing:7}
\emph{Coding.} Repeat the thousand-signal screen with null statistics that share a common factor (correlation 0.5). Does
Benjamini--Hochberg still control the \omterm{def:qm:testing-and-multiple-testing:fdr}{false discovery rate} at 5\%? How often does the realised proportion of false discoveries
exceed 10\%?
\end{exercise}

\begin{exercise}[$\star\star\star$]\label{exo:qm:testing-and-multiple-testing:8}
\emph{Find the flaw.} ``We tested our signal separately on twelve markets. It is significant at 5\% in three of them, where we now
trade it; the three backtests are our estimate of its performance.''
\end{exercise}

\section{Problem: Two Hundred Momentum Variants}

\begin{problem}[{Weekend problem --- how surprising is the best of a correlated search?}]\label{pb:qm:testing-and-multiple-testing:1}
A team backtests two hundred trend rules (hold the sign of the trailing $L$-day return, $L = 5, \dots, 204$) on ten years of a
market that, unknown to them, is a driftless random walk with 1\% daily volatility. The history is the chapter's seeded one.

\medskip\noindent\textbf{Part I --- The search.}
\begin{enumerate}
\item Which lookback wins, with what $t$-statistic and annualised \omterm{def:qm:estimation:sharpe}{Sharpe ratio}?
\item What is its naive one-sided $p$-value?
\item How many of the two hundred pass a one-sided 5\% test?
\item What is the mean $t$-statistic across the family, and why is it not near zero?
\item What is the correlation of the best rule's P\&L with its neighbour's, and with the 5-day and 204-day rules'?
\end{enumerate}

\medskip\noindent\textbf{Part II --- Corrections.}
\begin{enumerate}[resume]
\item What are the Bonferroni-adjusted $p$-value and the Bonferroni critical value?
\item What does Holm give for the best, and why the same?
\item What does Benjamini--Hochberg give, and why smaller?
\item What would the adjusted $p$-value be if the two hundred were independent?
\item What are the max-statistic adjusted $p$-value and critical value, from the estimated correlation?
\end{enumerate}

\medskip\noindent\textbf{Part III --- What the search amounts to.}
\begin{enumerate}[resume]
\item What is the effective number of independent trials?
\item How many variants clear the max-statistic threshold?
\item Over 2\,000 fresh histories, how often does the best rule reach $t = 3.11$?
\item What are the probabilistic \omterm{def:qm:estimation:sharpe}{Sharpe ratio} and the \omterm{def:qm:testing-and-multiple-testing:dsr}{deflated Sharpe ratio} with 200 and with the effective number of trials?
\item Do the answers change if daily returns have Student-$t$ tails with four degrees of freedom?
\end{enumerate}

\medskip\noindent\textbf{Part IV --- Judgement.}
\begin{enumerate}[resume]
\item The rules have no edge by construction. Why does the corrected test reject?
\item What would you require before trading the 42-day rule?
\item What should the research report record about the search?
\item State the \emph{named result}: the adjusted $p$-value of the best of two hundred correlated variants and the effective
number of trials.
\item In one sentence: what does a multiple-testing correction protect against, and what not?
\end{enumerate}
\end{problem}

\section{Interview questions}

\begin{interviewq}[$\star$ \iqroles{researcher, trader}]\label{iq:qm:testing-and-multiple-testing:1}
What is a $p$-value, and what is it not?
\end{interviewq}

\begin{interviewq}[$\star\star$ \iqroles{researcher, trader}]\label{iq:qm:testing-and-multiple-testing:2}
You backtested a hundred strategies and the best has $t = 2.8$. Is it real?
\end{interviewq}

\begin{interviewq}[$\star\star$ \iqroles{researcher, mle}]\label{iq:qm:testing-and-multiple-testing:3}
Bonferroni or Benjamini--Hochberg: when would you use which?
\end{interviewq}

\begin{interviewq}[$\star\star$ \iqroles{researcher, trader, risk}]\label{iq:qm:testing-and-multiple-testing:4}
How long must a strategy run before you can confirm a \omterm{def:qm:estimation:sharpe}{Sharpe ratio} of 1?
\end{interviewq}

\begin{interviewq}[$\star\star\star$ \iqroles{researcher}]\label{iq:qm:testing-and-multiple-testing:5}
Your two hundred variants are highly correlated. How do you correct for the search?
\end{interviewq}

\begin{interviewq}[$\star\star$ \iqroles{mle, researcher}]\label{iq:qm:testing-and-multiple-testing:6}
What is the \omterm{def:qm:testing-and-multiple-testing:snooping}{garden of forking paths}, and where does it hide in a machine-learning pipeline?
\end{interviewq}

\begin{interviewq}[$\star\star$ \iqroles{researcher, risk}]\label{iq:qm:testing-and-multiple-testing:7}
You fit a distribution to returns and a \omterm{def:qm:testing-and-multiple-testing:ks}{Kolmogorov--Smirnov test} does not reject it. What have you learned?
\end{interviewq}
